\subsection{Implementation Details}
The project utilizes several core classes divided into two main categories: Data Structures and UI Management.

\noindent\textbf{List of Core Classes and Relevant Functions:}
\begin{itemize}
    \item \textbf{\texttt{UI::Handler}:} The main engine.
    \begin{itemize}
        \item \texttt{run()}: Executes the main application loop, managing frames.
        \item \texttt{update\_camera()}: Handles user panning and zooming inputs.
    \end{itemize}
    \item \textbf{\texttt{UI::Canvas} (Polymorphic Base Class):} Defines the foundational virtual methods (\texttt{setup()}, \texttt{run()}, \texttt{update\_animation()}) to allow the \texttt{Handler} to manage different screens uniformly.
    \item \textbf{Derived Canvas Classes (e.g., \texttt{AVL\_Canvas}, \texttt{MST\_Canvas}):} Inherit from \texttt{UI::Canvas} to render specific structures. To maintain strict encapsulation, each derived class independently manages its own UI playback buttons (\texttt{next()}, \texttt{prev()}, \texttt{skip()}) and playback state variables (\texttt{current\_step}, \texttt{is\_playing}).
    \item \textbf{\texttt{Data\_Structure::[Type]} (e.g., \texttt{AVL\_Tree}, \texttt{Shortest\_Path}):}
    \begin{itemize}
        \item \texttt{insert()}, \texttt{erase()}, \texttt{find()}: Core algorithmic operations.
        \item \texttt{save\_snapshot()}: Captures the state of the structure at a specific step.
        \item \texttt{get\_copy()}: Creates a deep copy of the current memory state for history logging.
    \end{itemize}
\end{itemize}

\textbf{Main Flows and Class Collaboration:}
\begin{itemize}
    \item \textbf{Use Case 1: Data Insertion and History Logging Flow} \\
    When a user enters a value to insert into the AVL Tree, the \texttt{AVL\_Canvas} captures the input and parses it. It then calls \texttt{avl.insert(val)}. Inside the \texttt{AVL\_Tree} class, as the algorithm traverses the tree and performs rotations, it repeatedly calls \texttt{save\_snapshot(current\_line, OPERATION\_TYPE)}. This function uses \texttt{get\_copy()} to duplicate the tree's current pointers and nodes, pushing this frozen state into the \texttt{history} vector. This collaboration creates a chronological ledger of the algorithm's execution without disrupting the main data.

    \item \textbf{Use Case 2: Animation and Static Rendering Flow} \\
    Within the main game loop, the \texttt{UI::Handler} continuously invokes the active \texttt{Canvas}, meaning the \texttt{update\_animation()} method is actually called every single frame. However, its core interpolation logic is strictly guarded by an \texttt{is\_playing} boolean flag. If the user manually steps through the execution (e.g., clicking "Prev" or "Next" button), the Canvas retrieves the corresponding \texttt{Snapshot\_\allowbreak Data} and renders the static state. When \texttt{is\_playing} is true, the timer-based mechanism runs. The method computes the distance between current and target coordinates. If a node has not reached its destination, it smoothly interpolates the movement frame-by-frame. If it has already reached the target, it remains visually static while the timer continues to count down. Once the step duration expires, it automatically advances to the next historical snapshot, setting new target positions for the subsequent animation phase.
\end{itemize}